
\documentclass[11pt,a4paper]{article}
\usepackage[margin=2.2cm]{geometry}
\usepackage{enumitem}
\usepackage{tabularx}
\usepackage{booktabs}
\usepackage{hyperref}
\usepackage{parskip}
\hypersetup{colorlinks=true, linkcolor=blue, urlcolor=blue}
\newcommand{\notespace}{\vspace{0.3em}\noindent\rule{\linewidth}{0.4pt}\vspace{0.3em}}

\title{Geophysics IIIA: Week 2}
\author{Scenarios: Decision-Making for\\ Physical Property Mixing}
\date{}
\begin{document}
\maketitle
\noindent\textbf{Instructions:} Work in mixed geology/physics groups. For each scenario: (1) state the \emph{goal} and \emph{key properties}; (2) choose how you will \emph{obtain} those properties (measure, estimate from composition, infer from proxies); (3) select an \emph{appropriate mixing law} (justify); (4) identify \emph{uncertainties/non-uniqueness}; (5) draft \textbf{2--3 questions a geophysicist should ask a geologist} before committing to a survey/model. Use the note space provided.

\section*{Scenario 1 --- Gravity over Mixed Lithologies}
You are planning a gravity survey over shallow basement with weathered granite, mafic dyke swarms, and variably porous sediments.\newline
\textbf{Tasks:}
\begin{enumerate}[leftmargin=*]
  \item Which density components dominate the expected anomaly? When can density be predicted compositionally? When does porosity matter most?
  \item Choose a mixing law for sediments vs. crystalline units.
  \item Note assumptions and unit conventions.
  \item List major uncertainty sources and how you would reduce them.
  \item Write 2--3 questions for the geologist.
\end{enumerate}
\vspace{16pt}

\section*{Scenario 2 --- Magnetic Lineament Interpretation}
You have regional aeromagnetic data (20 km line spacing) over volcaniclastics and granitoids.\newline
\textbf{Tasks:}
\begin{enumerate}[leftmargin=*]
  \item Identify which property matters (susceptibility vs. remanence) and what fabric/mineralogy might control it.
  \item Would you estimate $\chi$ from lithology or require measurements? What are the risks of assuming induced-only magnetization?
  \item What geological information (e.g., thermal/alteration history) would change your interpretation most?
  \item Write 2--3 questions for the geologist.
\end{enumerate}
\vspace{16pt}

\section*{Scenario 3 --- Conductive Heat-Flow Study}
You need estimates of thermal conductivity $k(T)$ and heat production $A$ for three lithologies to model a geotherm.\newline
\textbf{Tasks:}
\begin{enumerate}[leftmargin=*]
  \item Select mixing laws for $k$ in random vs. layered media; justify.
  \item Decide how to estimate $A$ (e.g., from K, U, Th or lithology) and how uncertainty propagates into geotherms.
  \item Identify which additional measurements would reduce uncertainty most.
  \item Write 2--3 questions for the geologist.
\end{enumerate}
\vspace{16pt}

\section*{Scenario 4 --- Seismic Property Estimation}
You must propose a 1-D velocity model.\newline
\textbf{Tasks:}
\begin{enumerate}[leftmargin=*]
  \item Explain why Voigt--Reuss--Hill bounds are appropriate for moduli/velocities.
  \item Distinguish density effects from moduli effects; identify when cracks/porosity dominate.
  \item What rock fabric information would you request from geology to constrain anisotropy?
  \item Write 2--3 questions for the geologist.
\end{enumerate}
\vspace{16pt}

\section*{Score Sheet (Scenario Responses)}
\noindent\textbf{Student/Group:} \rule{0.42\linewidth}{0.4pt} \hfill \textbf{Scenario(s):} \rule{0.35\linewidth}{0.4pt}

\vspace{0.5em}
\begin{tabularx}{\linewidth}{@{}lXc@{}}
\toprule
Criterion & Description & Points \\
\midrule
Property strategy & Chooses appropriate properties and acquisition/estimation strategy with sound assumptions & 6 \\
Mixing-law choice & Selects and justifies a law suitable for geometry and property & 6 \\
Uncertainty & Identifies and prioritizes key uncertainties; proposes mitigation & 5 \\
Geo questions & High-value, concrete questions for the geologist & 5 \\
Clarity & Organized, legible, concise notes & 3 \\
\midrule
\textbf{Total} &  & \textbf{25} \\
\bottomrule
\end{tabularx}

\end{document}
